\section{Pesquisa Bibliográfica}

Como mencionado anteriormente, os trabalhos relacionados ao Mundo do Wumpus concentram-se, em sua maioria, no desenvolvimento e na avaliação de estratégias para agentes inteligentes, utilizando o ambiente como um cenário de testes para diferentes abordagens de Inteligência Artificial. Em geral, o foco dessas pesquisas está nos algoritmos e nos resultados obtidos pelos agentes, enquanto aspectos relacionados à acessibilidade, à integração das funcionalidades e ao suporte à experimentação por usuários finais recebem menor atenção.

Com o objetivo de identificar plataformas voltadas à execução e avaliação de agentes por meio de interfaces acessíveis ao usuário, foi realizada uma pesquisa bibliográfica abrangendo trabalhos que apresentassem simuladores ou ambientes de experimentação para o Mundo do Wumpus. Embora tenham sido encontradas diversas propostas de implementação do ambiente, observou-se que a maioria delas é destinada a fins específicos de pesquisa, frequentemente limitada a interfaces textuais ou à interação direta por meio de código. Além disso, as ferramentas que oferecem algum nível de interação gráfica ou configuração por parte do usuário correspondem, em sua maioria, a propostas desenvolvidas há mais de uma década e apresentam funcionalidades restritas quando comparadas às possibilidades oferecidas por plataformas web modernas.

Nesse contexto, a análise dos trabalhos encontrados é relevante para compreender como essas soluções foram concebidas, quais funcionalidades disponibilizam aos usuários e quais limitações ainda permanecem. Essa discussão também permite evidenciar a contribuição da plataforma proposta neste trabalho, que busca integrar, em um único ambiente web, recursos de criação de mundos, configuração de agentes, execução de experimentos e gerenciamento dos resultados obtidos.

\cite{julien2015wumpus} apresentam uma implementação do Mundo do Wumpus baseada em um sistema multiagente utilizando o framework JADE e programação lógica em Prolog. O trabalho concentra-se na modelagem do comportamento dos agentes, nos mecanismos de comunicação e nos algoritmos de inferência utilizados para a resolução do problema, disponibilizando uma interface gráfica para visualização do ambiente. Diferentemente da proposta apresentada neste trabalho, a ferramenta é concebida como um suporte à implementação e avaliação da arquitetura multiagente, não contemplando recursos voltados ao gerenciamento de experimentos e à interação ampliada do usuário.

\cite{salehi2012multicontext} propõem um ambiente de testes baseado em uma versão estendida do Mundo do Wumpus, projetada para simular restrições presentes em cenários reais de trabalho cooperativo entre agentes. A plataforma introduz novos elementos de complexidade, como múltiplos contextos, sensores imperfeitos e compartilhamento de informações entre agentes, sendo implementada em Java utilizando o framework JADE. Os autores destacam ainda a flexibilidade do ambiente e a disponibilidade de uma interface de configuração para apoiar a avaliação de problemas relacionados à cooperação, coordenação, planejamento e negociação em sistemas multiagentes. Embora a proposta apresente mecanismos de configuração e interação com o usuário, seu principal objetivo é servir como um ambiente especializado para a investigação de estratégias de cooperação em sistemas multiagentes.

\cite{li2025llmcave} apresentam o LLM-Cave, um ambiente baseado no Mundo do Wumpus desenvolvido para avaliação das capacidades de raciocínio e tomada de decisão de grandes modelos de linguagem. A plataforma adapta o problema clássico para uma interação totalmente textual, fornecendo ao agente descrições estruturadas do ambiente e métricas quantitativas de desempenho, como taxa de sucesso, recompensa média e custo computacional. O trabalho utiliza esse ambiente para investigar estratégias de inferência, como Chain of Speculation e Planner-Critic, aplicadas à navegação e resolução do problema. Embora o LLM-Cave forneça um ambiente flexível para a avaliação de agentes baseados em modelos de linguagem, seu foco está na análise do desempenho e das estratégias de inferência desses agentes.

\cite{lasisi2022alp4ai} apresentam o ALP4AI (Agent-based Learning Platform for Introductory Artificial Intelligence), uma plataforma educacional desenvolvida para apoiar o ensino introdutório de Inteligência Artificial. A ferramenta oferece uma interface gráfica para visualização do comportamento dos agentes e dos resultados obtidos, permitindo a configuração de diferentes ambientes baseados em grades bidimensionais, com suporte a cenários com obstáculos, múltiplos agentes e múltiplos objetivos. Além disso, a plataforma foi concebida para incentivar a aprendizagem prática, possibilitando que os estudantes implementem e avaliem algoritmos clássicos de busca em um ambiente visual e interativo. Embora o ALP4AI compartilhe com a presente proposta o objetivo de apoiar o ensino de Inteligência Artificial por meio de uma interface gráfica, sua aplicação está direcionada ao estudo de algoritmos de busca em um ambiente genérico baseado em estados.

\cite{bryce2011wumpus} apresenta uma abordagem para o ensino introdutório de Inteligência Artificial baseada em uma sequência de projetos desenvolvidos sobre o ambiente do Mundo do Wumpus. A proposta utiliza um único simulador como base para atividades práticas envolvendo diferentes tópicos da área, incluindo algoritmos de busca, inferência lógica, planejamento automatizado e redes bayesianas, permitindo que os estudantes concentrem seus esforços na compreensão dos conceitos de IA sem a necessidade de aprender múltiplos ambientes de desenvolvimento. O autor destaca ainda que a utilização de um domínio único favorece a aprendizagem prática (hands-on learning) e reduz a complexidade associada à adaptação a diferentes plataformas. Embora o objetivo do trabalho esteja voltado à organização de uma disciplina introdutória de IA, seus resultados evidenciam o potencial do Mundo do Wumpus como ferramenta de apoio ao ensino. Esse resultado reforça a importância do desenvolvimento de ferramentas que tornem a experimentação no Mundo do Wumpus mais acessível, favorecendo tanto atividades de ensino quanto aplicações de pesquisa.

De maneira geral, os trabalhos analisados evidenciam a relevância do Mundo do Wumpus como ambiente para avaliação de agentes inteligentes e para o ensino de conceitos fundamentais de Inteligência Artificial. Entretanto, observa-se que a maior parte das propostas encontradas está voltada a objetivos específicos, como a investigação de arquiteturas multiagentes, estratégias de cooperação, avaliação de modelos de linguagem ou ensino de algoritmos clássicos, oferecendo recursos limitados para criação e gerenciamento de experimentos por meio de uma plataforma integrada e acessível. Nesse contexto, o protótipo apresentado neste trabalho busca complementar essas iniciativas ao reunir, em um único ambiente web, funcionalidades para criação de mundos, configuração de agentes, execução de simulações, armazenamento de resultados e análise de desempenho, atendendo tanto a aplicações educacionais quanto a atividades de pesquisa.
 